\documentclass[11pt]{article}
\usepackage{amsmath,amssymb,amsthm,booktabs,hyperref}
\newtheorem{theorem}{Theorem}
\newtheorem{lemma}[theorem]{Lemma}
\newtheorem{corollary}[theorem]{Corollary}
\title{Incidence matrices for graphs with two isolated middle vertices}
\author{human\_compression project}
\date{}
\begin{document}\maketitle

\paragraph{Scope and attribution.}
The problem is the ordinary, unrooted graph count for order $32$, degree
$5$, girth $4$ and diameter $8$ in Cambie, Goedgebeur, Jooken and
Van den Eede, \emph{On the order-diameter ratio of girth-diameter cages},
DMTCS 28:4, SOFSEM 2026 \#1, arXiv:2511.21144v3, 19 August 2026,
Appendix B, Table 6, physical page 23. Its final parameter is diameter.
No edge-girth-regular multiplicity condition is assumed.

The independent-neighborhood/defect setup and side-incidence viewpoint
are borrowed from the manuscript and the retained F009 reduction notes.
The present assembly consists of elementary degree, vertex-cover,
false-twin and distance arguments. No priority or novelty is claimed.
This note proves a conditional template theorem; it neither matches a
pending search family nor increases the accepted census count. In
particular, the exact two-isolated-middle result below is an exhaustive
parameterization, not an asserted number of ordinary isomorphism classes.
No Source, graph representation, field, generator, state or checker was
executed in preparing it.

\paragraph{Ambient definitions.}
Let $G$ be connected, simple, triangle-free, $5$-regular, on $32$
vertices, of diameter eight. Fix a diametral geodesic $p_0,\ldots,p_8$.
Put
\[
 A=N(p_0),\quad B=N(p_4),\quad C=N(p_8),\quad I=A\cup B\cup C,
 \quad J=V(G)\setminus I,\quad F=G[J].
\]
On $K=J\setminus\{p_0,p_4,p_8\}$ put
\[
 L=\{x:N(x)\cap A\ne\varnothing\},\quad
 R=\{x:N(x)\cap C\ne\varnothing\},\quad M=K\setminus(L\cup R).
\]
All definitions refer to a complete graph. The profile considered here
is $M$ consisting of exactly $q$ vertices, all isolated in $F$.
There may not be an additional nonisolated vertex of $M$.

\begin{lemma}[Basic side restrictions]\label{basic}
The set $I$ is independent of size $15$, $|J|=17$, and $e(F)=5$.
The centers are isolated in $F$, the sides $L,R$ are disjoint, neither
has fewer than five vertices, and no $F$-edge joins them. A vertex of
$B$ cannot have a neighbor in both sides. For $D=G[L]$,
$d_D(x)\leq |L|-4$, with the analogous statement on $R$.
If all $q$ middle vertices are isolated, then
\[
 H=\{p_4\}\cup M\quad\hbox{has identical open neighborhood }B,
 \qquad q\leq3.
\]
Every $b\in B$ has precisely $r=4-q$ neighbors, all in just one
side. Thus $B=B_L\sqcup B_R$, where both parts are nonempty.
\end{lemma}
\begin{proof}
The centers are at mutual distances at least four. Their neighborhoods
are pairwise disjoint and have no edges within or between them, by
triangle-freeness and the center distances. The $80$ edges of $G$
include $75$ incidences from $I$ to $J$, leaving five in $F$.
A vertex meeting both $A$ and $C$ shortens the diametral path to
length four; an $L$--$R$ edge shortens it to five. A $B$-vertex
meeting both sides gives a path of length six from $p_0$ to $p_8$.
These are impossible.

Each $A$-vertex needs four distinct $L$-neighbors besides $p_0$.
If $|L|=4$, all of them have all five neighbors in $A$, giving a
proper $K_{5,5}$ component. Hence $|L|\geq5$. Given $x\in L$,
choose a neighboring $A$-vertex. Its other three side neighbors
are nonneighbors of $x$, so $d_D(x)\leq |L|-4$.
An isolated $h\in M$ can meet only $B$, and degree five gives
$N(h)=B$. Four such vertices together with $p_4$ would close another
$K_{5,5}$ component. This proves $q\leq3$ and the port count.
If all ports went to one side, the other side and its endpoint
neighborhood and center would be a proper component.
\end{proof}

\begin{theorem}[Exact side-incidence description]\label{hyper}
Suppose $M$ consists of exactly $q\in\{1,2,3\}$ isolated vertices.
The standalone side-data conditions explicitly include the bounds of
Lemma~\ref{basic}: $n_L,n_R\geq5$ and $1\leq b_L,b_R\leq4$.
In particular both port parts are nonempty.
A side with $n$ vertices and $b$ ports is described by:
\begin{enumerate}
\item a triangle-free simple graph $D$ on $n$ side vertices;
\item a multiset $\mathcal O$ of five $(n-4)$-subsets, each a
vertex cover of $D$;
\item a multiset $\mathcal P$ of $b$ independent $r$-subsets of $D$,
where $r=4-q$.
\end{enumerate}
Write $o(x),k(x)$ for the respective occurrence counts. The exact
degree and side-membership conditions are
\[
 o(x)-k(x)=d_D(x),\qquad o(x)\leq4.                 \tag{1}
\]
In addition,
\[
 5n=20+rb+2e(D),\quad \Delta(D)\leq n-4.           \tag{2}
\]
The two sides have $n_L+n_R=14-q$, $b_L+b_R=5$ and
$e(D_L)+e(D_R)=5$.

Conversely, two such side descriptions construct a unique graph up to
the indicated vertex relabelings: take an independent set $H$ of size
$q+1$, all five ports adjacent to all of $H$, and on each side add an
endpoint center and five independent endpoint-neighborhood vertices.
The latter all meet that center; the row corresponding to $O\in
\mathcal O$ meets precisely $V(D)\setminus O$. Each port meets its
listed $P\in\mathcal P$. Include the edges of $D$ and no others.
This graph is simple, triangle-free and $5$-regular on $32$ vertices.
If it has diameter eight, its designated centers recover exactly the
stated profile. In the $q=2$ case, diameter eight follows automatically
from the subsequent theorem.
\end{theorem}
\begin{proof}
An endpoint-neighborhood row has four side neighbors. Their complement
has size $n-4$ and covers every internal edge, exactly because the four
neighbors are independent. A port also has an independent neighbor set,
by triangle-freeness. The degree at $x$ is
$5-o(x)+k(x)+d_D(x)$, giving (1), while $o(x)\leq4$ is exactly the
condition that $x$ has an endpoint-neighborhood neighbor. Summing the
degree equation proves (2).

For the converse, endpoint rows have degree five, ports have degree
$(q+1)+r=5$, vertices in $H$ and the endpoint centers have degree five,
and (1) supplies the other degrees. The independent-row conditions and
triangle-freeness of $D$ exclude every possible triangle. All side
vertices meet their endpoint neighborhood, all ports meet their side,
and both port parts are nonempty, so the graph is connected.
There are $14-q$ side vertices, $q+1$ middle-center vertices, two
endpoint centers, ten endpoint-neighborhood vertices and five ports:
the total is $32$. Every side vertex meets its designated endpoint
neighborhood, whereas precisely the other $q$ vertices of $H$ meet
neither. Thus the profile is exact whenever the centers lie on a
diametral path as stipulated.
\end{proof}

\begin{theorem}[Two isolated middle vertices: exhaustive small templates]
\label{two}
Suppose $M$ consists of exactly two isolated vertices. Both sides have
six vertices. Each side is completely described by a triple $(D,U,V)$,
where $U$ and $V$ are loopless multigraphs on these six vertices:
\[
 e(U)=5,\quad e(V)=b,\quad 1\leq b\leq4,\quad e(D)=5-b;
                                                               \tag{3}
\]
every edge of $U$ is a two-vertex cover of $D$; every edge of $V$
is an independent pair in $D$; and
\[
 d_U(x)-d_V(x)=d_D(x),\qquad d_U(x)\leq4.          \tag{4}
\]
Parallel edges are allowed in $U,V$; loops are not. The possible
side shapes are exactly the shapes not excluded in this table, with
attachment multiplicities still subject to (4):
\[
\begin{array}{c|c|l}
b&e(D)&D\\\hline
1&4&C_4\sqcup2K_1\\
2&3&P_4\sqcup2K_1\\
3&2&P_3\sqcup3K_1\ \hbox{or}\ 2P_2\sqcup2K_1\\
4&1&P_2\sqcup4K_1.
\end{array}                                                    \tag{5}
\]
Pair any two admissible side triples with port counts $(1,4)$ or
$(2,3)$, allowing side interchange, and use Theorem~\ref{hyper} with
$H$ of size three. Every resulting graph has girth four and diameter
eight, and is nonbipartite. Every graph in the stated profile occurs.
Thus (3)--(5), or the explicit parameter forms below, are a complete
finite template description; no additional graph-search predicate
is needed for degree, triangle-freeness, connectedness or diameter.
\end{theorem}
\begin{proof}[Sizes and shapes]
The side sizes sum to twelve and are at least five. If one had size
five, its degree sum would be $25=20+2b+2e(D)$, impossible by parity.
Hence both sizes are six. Equation (2) becomes $e(D)+b=5$, and the
maximum side degree is two. The five edges of $F$ are precisely the
internal side edges.

For four edges on six vertices, the triangle-free maximum-degree-two
possibilities are $C_4\sqcup2K_1$, $P_5\sqcup K_1$,
$P_4\sqcup P_2$ and $2P_3$. The third has no independent four-set.
In the second and fourth every independent four-set contains a fixed
nonisolated vertex; that vertex would have all five endpoint-row
neighbors as well as positive internal degree. They are impossible.
For three edges the possibilities are $P_4\sqcup2K_1$,
$P_3\sqcup P_2\sqcup K_1$ and $3P_2$. The last has independence
number three; in the middle every independent four-set contains both
endpoints of the $P_3$, again violating degree five. Two edges give
exactly the two listed shapes, and one edge gives the final shape.
This proves exhaustion of (5); conditions (4) decide the remaining
attachments and their existence.
\end{proof}

\begin{proof}[Automatic diameter and exact profile]
Fix a side and let $S$ be the union of its port-pair supports.
It has at least two vertices. Each endpoint row has four side
neighbors. If this four-set meets $S$, the row has distance three
to any vertex of $H$. Otherwise $S$ has exactly two vertices and
is its omitted pair, hence is a vertex cover of $D$. Since $e(D)>0$
and $S$ is independent, some $D$-edge runs from that four-set to
$S$. The row then has distance four to $H$.

A side vertex $x\in S$ has distance two to $H$. Suppose $x\notin S$.
If an endpoint row contains $x$ and some element of $S$, it gives
an $x$--$H$ path of length four. Otherwise every row containing $x$
omits all of $S$. Such a row exists by $d_U(x)\leq4$, so $|S|=2$
and $S$ covers $D$. If $d_D(x)>0$, an internal neighbor lies in
$S$, yielding distance three to $H$. If $d_D(x)=0$, then
$d_V(x)=0$ and (4) gives $d_U(x)=0$. Thus all five endpoint rows
contain $x$, and under the supposition all five omit $S$. But
vertices of $S$ are also side vertices and must meet an endpoint
row, a contradiction.

Every vertex on a side therefore has distance at most four to $H$;
its endpoint center has distance exactly four. Indeed a port support
vertex has an endpoint-row neighbor, giving a path of length four,
and the prescribed construction has no shorter center-to-$H$ path.
Deleting $H$ separates the sides, so the endpoint centers are at
distance eight, and all graph distances are at most eight through
$H$. The three chosen centers occur at positions zero, four and
eight of a shortest path. The other two vertices of $H$ are exactly
the isolated middle vertices. A four-cycle uses two vertices of $H$
and two ports. Since every internal side edge has endpoints with
nonempty, disjoint endpoint-row neighborhoods, it also gives a
five-cycle through the endpoint center. Thus girth is exactly four
and the graph is nonbipartite.
\end{proof}

\paragraph{Explicit attachment parameters: one-port side.}
For $D=C_4\sqcup2K_1$, let $P,Q$ be the opposite pairs of the
four-cycle. The only two-covers are $P,Q$. Their multiplicities in
$U$ are two and three, up to interchange. The sole $V$-edge is
the pair having multiplicity three. This is the unique one-port
side type with the endpoint center distinguished. The endpoint
center and the two isolates are false twins; each opposite pair
of the cycle also consists of false twins. Its weighted quotient
has six vertices: the endpoint class of weight three, the two
cycle-pair classes of weight two, two endpoint-row classes of
weights two and three, and the single port.

\paragraph{Explicit attachment parameters: two-port side.}
Label the path $1$--$2$--$3$--$4$, with isolates $5,6$.
The possible $U$-edges are $13,23,24$, of multiplicities
$a,5-a-c,c$. The only $V$-edges that can occur are $13,14,24$,
with multiplicities
\[
 3-c,\quad a+c-4,\quad 3-a,\qquad
 1\leq a,c\leq3,\quad4\leq a+c\leq5.              \tag{6}
\]
These formulas follow by solving the four vertex-degree equations;
the isolates have degree zero in both multigraphs. Reflection of
the path exchanges $a,c$. There are exactly three rooted side
types, represented by $(a,c)=(1,3),(2,2),(2,3)$. The center and
the two isolates are false twins. There are at most three
endpoint-row classes and two port classes, so the full side
false-twin quotient has at most ten vertices.

\paragraph{Explicit attachment parameters: three-port side, matching.}
For $D$ with edges $12,34$ and isolates $5,6$, both $U$ and $V$
use only cross-pairs between $\{1,2\}$ and $\{3,4\}$. Regard their
multiplicities as nonnegative $2\times2$ integer matrices. All
solutions are exactly
\[
 \sum U_{ij}=5,\quad \sum V_{ij}=3,\qquad
 U\mathbf1=V\mathbf1+\mathbf1,\quad
 U^T\mathbf1=V^T\mathbf1+\mathbf1.                 \tag{7}
\]
The margin conditions make every degree of $U$ at most four.
Row swaps, column swaps and transpose act simultaneously on the
two matrices, and describe precisely the rooted side relabelings.
The isolates join the endpoint center as a false-twin class.

\paragraph{Explicit attachment parameters: three-port side, path.}
Let $D$ be the path $u$--$v$--$w$ with isolates $i,j,k$.
The $U$-edges are $uw$ of multiplicity $t$, $vu$ of multiplicity
$a$, $vw$ of multiplicity $c$, and $v\ell$ of multiplicities
$z_\ell$ for $\ell\in\{i,j,k\}$. They obey
\[
 1\leq t\leq3,\quad a,c,z_\ell\geq0,\quad
 t+a+c+\sum z_\ell=5,\quad t+a,t+c\leq4.           \tag{8}
\]
Choose any loopless three-edge multigraph $V$ using only
$uw$, $u\ell$, $w\ell$, $v\ell$ and $\ell\ell'$ for distinct
isolates, with degree sequence
\[
 d_V(u)=t+a-1,\quad d_V(w)=t+c-1,\quad
 d_V(v)=3-t,\quad d_V(\ell)=z_\ell.                \tag{9}
\]
Equations (8)--(9) are necessary and sufficient; no unlisted
edge is allowed. The lower bound $t\geq1$ follows because
$d_U(v)=5-t\leq4$; the upper bound is $d_V(v)\geq0$.
The rooted equivalences are $u\leftrightarrow w$ and permutations
of the isolates, applied to every parameter and edge together.
No numerical count of these side types is needed here.

\paragraph{Explicit attachment parameters: four-port side.}
Let $D$ have its sole edge $uv$ and isolates $1,2,3,4$.
Write the $U$-multiplicities as
\[
 U_{uv}=x,\quad U_{u i}=a_i,\quad U_{v i}=c_i,
 \qquad x+\sum_i(a_i+c_i)=5.                        \tag{10}
\]
Put $A=\sum a_i$, $C=\sum c_i$.
All solutions have \emph{either} $x=1$ \emph{or} $x=0$:
\begin{itemize}
\item If $x=1$, choose nonnegative $s_i,t_i$ with
\[
 A+C=4,\quad1\leq A,C\leq3,\quad
 \sum s_i=A,\quad\sum t_i=C,\quad
 s_i+t_i=a_i+c_i.                                 \tag{11}
\]
The $V$-edges are $u i$ of multiplicity $s_i$ and $v i$ of
multiplicity $t_i$.
\item If $x=0$, choose a distinguished unordered pair $\{j,k\}$
of distinct isolates and nonnegative $s_i,t_i$ with
\[
 A+C=5,\quad1\leq A,C\leq4,\quad a_i+c_i\leq4,
 \quad\sum s_i=A-1,\quad\sum t_i=C-1,
 \quad s_i+t_i=a_i+c_i-\mathbf1_{i\in\{j,k\}}.      \tag{12}
\]
The $V$-edges are $u i$ with multiplicity $s_i$, $v i$ with
multiplicity $t_i$, and exactly one edge $jk$.
\end{itemize}
To prove completeness of this split, all $V$-neighbors of $u,v$
are isolates. Their total required degree is $3+x$, whereas the
total degree required on the isolates is $5-x$. If $z$ is the
number of isolate--isolate $V$-edges, this gives $z=1-x$.
Hence $x\in\{0,1\}$ and the two forms follow. Conversely their
degree equations, permitted edges and upper bounds give exactly
(3)--(4). Relabeling is the simultaneous interchange of $u,v$
and arbitrary permutation of the four isolates.

\begin{theorem}[Ordinary isomorphism by the full false-twin quotient]
\label{twins}
For any finite simple graph, identify vertices with identical open
neighborhoods, retain their class sizes as positive integer weights,
and retain adjacency of classes. The graph is exactly the independent
blow-up of this weighted quotient. Two graphs are ordinarily
isomorphic if and only if these full weighted quotients are isomorphic.
The quotient is intrinsic; it uses no root, side or row colors.

Consequently the two-isolated-middle templates above admit a complete
ordinary-isomorphism test by this quotient, including coincidences
between different root choices or port-count decompositions. Every
such quotient has at most $26$ vertices. This is an upper bound,
not a claimed number of ordinary classes or a group-order quotient.
\end{theorem}
\begin{proof}
Distinct vertices with equal open neighborhoods are nonadjacent.
Adjacency between their classes is complete or empty. Therefore
the original graph is exactly recovered by the weighted expansion.
An isomorphism preserves equality of neighborhoods and class sizes,
and hence induces a weighted quotient isomorphism. Conversely a
weighted quotient isomorphism can be lifted by arbitrary bijections
of equally sized classes.

For the size bound, $H$ of size three always removes two quotient
vertices. A one-port side has $13$ original vertices and quotient
size six, as described above. Its complementary four-port side
has $16$ vertices, so a $(1,4)$ pair has quotient size at most
$1+6+16=23$. A two-port side has $14$ original vertices and at
most ten quotient vertices. A three-port side has $15$ original
vertices. Thus a $(2,3)$ pair has quotient size at most
$1+10+15=26$. Any additional false twins only improve these bounds.
The use of the \emph{full} quotient is essential: a partition by
selected row labels or guessed automorphism orbits is insufficient.
\end{proof}

\paragraph{An explicit two-isolated template beyond the three-isolated types.}
Take a two-port path side $1$--$2$--$3$--$4$, with two isolates,
and
\[
 U_2=2[13]+[23]+2[24],\qquad V_2=[13]+[24].
\]
Take a three-port path side $u$--$v$--$w$, with isolates $i,j,k$,
and
\[
 U_3=[uw]+[vu]+[vi]+[vj]+[vk],\qquad
 V_3=[uk]+[vi]+[vj].                               \tag{13}
\]
Bracket notation denotes a multigraph edge with the displayed
multiplicity. Equations (6), (8) and (9) hold, so the preceding
proof constructs a graph of the required order, degree, girth,
diameter and exact two-isolated profile without a search.

Its full false-twin quotient has exactly $26$ vertices. On the
two-port side the center and two isolates form one weight-three
class, the three distinct endpoint rows have weights $2,1,2$,
the four path vertices are distinct classes, and the two ports
are distinct classes: ten in all. On the three-port side all
fifteen vertices have different open neighborhoods: the five
endpoint rows and three ports have different supports; the
internal-path neighbor sets distinguish $v$ from $u,w$, and the
different endpoint-row degrees distinguish $u,w$; each isolate
has its own omitted row and its own port, distinguishing them
from each other and the center. The middle class has weight
three. Classes on opposite sides cannot coincide, since each
side vertex meets its own endpoint neighborhood and every
endpoint row contains its own center; ports contain $H$.
No endpoint-row/port class can coincide with a $J$-class: the
former have all neighbors in $J$, whereas every $J$-vertex has
at least one neighbor in $I$. Thus $10+15+1=26$.

The separately proved exact-three-isolated-middle classification
has full quotient sizes $15$ and $16$. Hence (13) is ordinarily
nonisomorphic to either of those two representatives, regardless
of possible alternative spine choices. This is a structural
existence and inequivalence statement, not a newly discovered
census class or an assertion that an outstanding family contains it.

\paragraph{What this can replace, and what remains owed.}
For an exact $M=2$ isolated profile, the entire graph is replaced by
two six-vertex side shapes and the small nonnegative integer tables
(6)--(12), with row repetitions represented as weights. Verifying one
template requires only its explicit incidence, degree and full
false-twin expansion; its girth and diameter follow from the proof.
There is no need to enumerate the full automorphism group or the
ambient space of all $32$-vertex graphs.

An application to a partial worker family must independently prove
that \emph{every} completion has exactly this profile, and must retain
all original required/forbidden edges, spine roles, label transports,
canonical histories and unvisited suffixes. A template is usable only
if some complete expansion satisfies all those constraints; a family
may realize no template or several. Distinct families or different
spines may realize the same ordinary graph. The intrinsic weighted
quotient proves how to identify such coincidences, but this note has
not run that identification or counted its equivalence classes.
No pending family, checkpoint, manuscript, reference total or lower
bound is paid merely by this conditional theorem.
\end{document}
